% Generated by officeParser.
\RequirePackage{iftex}
% DVI output (latex, uplatex or platex, then dvipdfmx): a global class option gives every package the dvipdfmx driver.
\def\officeparserdriver{}\ifpdf\else\ifXeTeX\else\def\officeparserdriver{dvipdfmx,}\fi\fi
\documentclass[\officeparserdriver 12pt]{article}
\iftutex
  \usepackage{fontspec}
\else
  \usepackage[T1]{fontenc}
  \usepackage[utf8]{inputenc}
  \usepackage{lmodern}
\fi
\usepackage[paperwidth=595.28pt,paperheight=841.89pt,top=72pt,bottom=72pt,left=72pt,right=72pt]{geometry}
\usepackage{parskip}
\usepackage{graphicx}
\usepackage{array,longtable}
\usepackage[normalem]{ulem}
\usepackage{xcolor}
\usepackage[hyperfootnotes=false]{hyperref}
\hypersetup{colorlinks=true,
  linkcolor=blue,
  urlcolor=blue,
  citecolor=blue,
  pdfcreationdate={D:20251128153644Z},
  pdfmoddate={D:20251226101743Z},
  pdfinfo={TestString={Hello from custom props},TestNumber={42},TestBool={true}}}
\setcounter{secnumdepth}{-\maxdimen}
\definecolor{hex000000}{HTML}{000000}
\definecolor{hex0000FF}{HTML}{0000FF}
\definecolor{hex0070C0}{HTML}{0070C0}
\definecolor{hex4F81BD}{HTML}{4F81BD}
\definecolor{hex808080}{HTML}{808080}
\definecolor{hex92D050}{HTML}{92D050}
\definecolor{hexFF0000}{HTML}{FF0000}

\begin{document}

\section{Sheet1}

\begin{longtable}{|*{2}{>{\raggedright\arraybackslash}p{\dimexpr(\linewidth-4\tabcolsep-3\arrayrulewidth)/2\relax}|}}
\hline
Demonstration of DOCX support in~calibre~ &  \\
\hline
\textcolor{hex000000}{calibre} &  \\
\hline
There is support for images, tables, lists, footnotes, endnotes,~links,~dropcaps~and~various~types~of text and paragraph level formatting.~ &  \\
\hline
\textcolor{hex000000}{\textbf{“Add Books”~Convert”.~“OK”}} &  \\
\hline
 &  \\
\hline
 &  \\
\hline
 &  \\
\hline
 &  \\
\hline
 &  \\
\hline
 &  \\
\hline
 &  \\
\hline
 &  \\
\hline
 &  \\
\hline
 &  \\
\hline
 &  \\
\hline
 &  \\
\hline
 &  \\
\hline
 &  \\
\hline
 &  \\
\hline
 &  \\
\hline
 &  \\
\hline
 &  \\
\hline
 &  \\
\hline
 &  \\
\hline
 &  \\
\hline
 &  \\
\hline
 &  \\
\hline
 &  \\
\hline
 &  \\
\hline
 &  \\
\hline
 &  \\
\hline
 &  \\
\hline
 &  \\
\hline
 &  \\
\hline
 &  \\
\hline
 &  \\
\hline
 &  \\
\hline
 &  \\
\hline
 &  \\
\hline
 &  \\
\hline
 &  \\
\hline
 &  \\
\hline
 &  \\
\hline
 &  \\
\hline
 &  \\
\hline
 &  \\
\hline
 &  \\
\hline
 &  \\
\hline
 &  \\
\hline
 &  \\
\hline
 &  \\
\hline
 &  \\
\hline
 &  \\
\hline
 &  \\
\hline
 &  \\
\hline
 &  \\
\hline
 &  \\
\hline
 &  \\
\hline
 &  \\
\hline
 &  \\
\hline
 &  \\
\hline
 &  \\
\hline
 &  \\
\hline
 &  \\
\hline
 &  \\
\hline
 &  \\
\hline
 &  \\
\hline
 &  \\
\hline
 &  \\
\hline
 &  \\
\hline
 &  \\
\hline
 &  \\
\hline
 &  \\
\hline
 &  \\
\hline
 &  \\
\hline
 &  \\
\hline
 &  \\
\hline
 &  \\
\hline
 &  \\
\hline
 &  \\
\hline
 &  \\
\hline
 &  \\
\hline
 &  \\
\hline
 &  \\
\hline
 &  \\
\hline
 &  \\
\hline
 &  \\
\hline
 &  \\
\hline
 &  \\
\hline
 &  \\
\hline
 &  \\
\hline
 &  \\
\hline
 &  \\
\hline
 &  \\
\hline
 &  \\
\hline
 &  \\
\hline
 &  \\
\hline
 &  \\
\hline
 &  \\
\hline
 &  \\
\hline
 &  \\
\hline
 &  \\
\hline
 &  \\
\hline
 &  \\
\hline
 &  \\
\hline
 &  \\
\hline
 &  \\
\hline
 &  \\
\hline
 &  \\
\hline
 &  \\
\hline
 &  \\
\hline
 &  \\
\hline
 &  \\
\hline
 &  \\
\hline
 &  \\
\hline
 &  \\
\hline
 &  \\
\hline
 &  \\
\hline
 &  \\
\hline
 &  \\
\hline
 &  \\
\hline
 & One, with~a very long~line to~demonstrate~that the hanging indent for the list is working correctly~ \\
\hline
 & Two~ \\
\hline
 &  \\
\hline
 & One~ \\
\hline
 & Two~ \\
\hline
 & Three~ \\
\hline
 & Four with~a very long~line to~demonstrate~that the hanging indent for the list is working correctly.~ \\
\hline
 & Five~ \\
\hline
 & Six~ \\
\hline
 &  \\
\hline
 &  \\
\hline
 &  \\
\hline
 &  \\
\hline
 &  \\
\hline
 &  \\
\hline
 &  \\
\hline
 & One~ \\
\hline
 & Two~ \\
\hline
 &  \\
\hline
 & We now resume our normal programming~ \\
\hline
 & Four~ \\
\hline
\end{longtable}

\section{Sheet2}

\begin{longtable}{|*{1}{>{\raggedright\arraybackslash}p{\dimexpr(\linewidth-2\tabcolsep-2\arrayrulewidth)/1\relax}|}}
\hline
Text Formatting~ \\
\hline
Inline formatting~ \\
\hline
Here, we~demonstrate~various types~of inline text formatting and the use of embedded fonts.~ \\
\hline
\textcolor{hex000000}{Here~is~some~}\textcolor{hex000000}{\textbf{bold,~}}\textcolor{hex000000}{\textit{italic,~}}\textcolor{hex000000}{\textit{\textbf{bold-italic,~}}}\textcolor{hex000000}{\uline{underlined~}}\textcolor{hex000000}{and~}\textcolor{hex000000}{\sout{struck~out~}}\textcolor{hex000000}{~text. Then, we have a super}\textcolor{hex000000}{{\fontsize{9pt}{10.8pt}\selectfont script}}\textcolor{hex000000}{~and a sub}\textcolor{hex000000}{{\fontsize{9pt}{10.8pt}\selectfont script}}\textcolor{hex000000}{. Now we see some~}\textcolor{hexFF0000}{red}\textcolor{hex000000}{,~}\textcolor{hex92D050}{green}\textcolor{hex000000}{~and~}\textcolor{hex0070C0}{blue}\textcolor{hex000000}{~text. Some text with a~yellow highlight. Some text in a~box. Some text~in~} \\
\hline
inverse video \\
\hline
. \\
\hline
\textcolor{hex808080}{\textit{subtle~emphasis~~}}\textcolor{hex000000}{\textbf{strong text~}}\textcolor{hex4F81BD}{\textit{\textbf{intense emphasis}}} \\
\hline
Fun with fonts~ \\
\hline
\textcolor{hex000000}{\texttt{some text in the Ubuntu Mono typeface, notice how every letter has the same width, even~i~and m}} \\
\hline
Paragraph level formatting~ \\
\hline
You can do crazy things with~paragraphs, if~the urge strikes you. For~instance~this paragraph is~right aligned~and has a right border. It has also been given a light gray background.~ \\
\hline
For~the lovers~of poetry amongst you, paragraphs with hanging indents, like this oftencome in handy. You can use hanging indents to ensure that a line of poetry~retains~its individual identity as a line even when the screen~is~ too~narrow to display it as a single line. Not only does this paragraph have a hanging indent,~it~is also has an extra top margin, setting it apart from the preceding paragraph.~ \\
\hline
\end{longtable}

\section{Sheet3}

{\small\setlength{\tabcolsep}{4pt}
\begin{longtable}{|*{6}{>{\raggedright\arraybackslash}p{\dimexpr(\linewidth-12\tabcolsep-7\arrayrulewidth)/6\relax}|}}
\hline
Tables~ &  &  &  &  &  \\
\hline
ITEM~ & NEEDED~ &  &  &  &  \\
\hline
Books~ & 1~ &  &  &  &  \\
\hline
Pens~ & 3~ &  &  &  &  \\
\hline
Pencils~ & 2~ &  &  &  &  \\
\hline
Highlighter~ & 2 colors~ &  &  &  &  \\
\hline
Scissors~ & 1 pair~ &  &  &  &  \\
\hline
 &  &  &  &  &  \\
\hline
Tables in Word can vary from the extremely simple to the extremely complex.~calibre~tries to do~its~best when converting tables. While you may run into trouble with the occasional table,~the vast majority of~common cases should be converted very well, as~demonstrated~in this section.~Note that for~optimum~results, when creating tables in Word, you should set their widths using percentages, rather than absolute units.~~To the left of this paragraph is a floating two column table with a nice green border and header row.~ &  &  &  &  &  \\
\hline
Now~let’s~look at a fancier table—one with alternating row colors and partial borders. This table is stretched out to take 100\% of the available width.~ &  &  &  &  &  \\
\hline
City or Town~ & Point A~ & Point B~ & Point C~ & Point D~ & Point E~ \\
\hline
Point A~ & —~ &  &  &  &  \\
\hline
Point B~ & 87~ & —~ &  &  &  \\
\hline
Point C~ & 64~ & 56~ & —~ &  &  \\
\hline
Point D~ & 37~ & 32~ & 91~ & —~ &  \\
\hline
Point E~ & 93~ & 35~ & 54~ & 43~ & —~ \\
\hline
\end{longtable}
}

\section{Sheet4}

\begin{longtable}{|*{4}{>{\raggedright\arraybackslash}p{\dimexpr(\linewidth-8\tabcolsep-5\arrayrulewidth)/4\relax}|}}
\hline
Next, we see a table with special formatting in various locations. Notice how the formatting for the header row and sub header rows is preserved.~ &  &  &  \\
\hline
 &  &  &  \\
\hline
College~ & New students~ & Graduating students~ & Change~ \\
\hline
 & ~ &  &  \\
\hline
Cedar University~ & 110~ & 103~ & +7~ \\
\hline
Oak Institute~ & 202~ & 210~ & -8~ \\
\hline
 & ~ &  &  \\
\hline
Cedar University~ & 24~ & 20~ & +4~ \\
\hline
Elm College~ & 43~ & 53~ & -10~ \\
\hline
Total~ & 998~ & 908~ & 90~ \\
\hline
 &  &  &  \\
\hline
 &  &  &  \\
\hline
\textcolor{hex000000}{{\fontsize{10pt}{12pt}\selectfont ~Fictitious data, for illustration purposes only~}} &  &  &  \\
\hline
 &  &  &  \\
\hline
 &  &  &  \\
\hline
 &  &  &  \\
\hline
Next, we have something a little more complex, a nested table,~i.e.~a table inside another table. Additionally, the inner table has some of its cells merged.~The table is displayed horizontally centered.~ &  &  &  \\
\hline
 &  &  &  \\
\hline
One~ & Two~ & To the left is a table inside a table, with some cells merged.~ &  \\
\hline
Three~ &  &  &  \\
\hline
 & Four~ &  &  \\
\hline
 &  &  &  \\
\hline
\end{longtable}

\section{Sheet5}

{\scriptsize\setlength{\tabcolsep}{2pt}
\begin{longtable}{|*{13}{>{\raggedright\arraybackslash}p{\dimexpr(\linewidth-26\tabcolsep-14\arrayrulewidth)/13\relax}|}}
\hline
We~end~with a~fancy~calendar, note how~much of the original formatting is preserved.~Note that this table will only display correctly on~relatively wide~screens. In general, very wide tables or tables whose cells have fixed width requirements~don’t~fare well in~ebooks.~ &  &  &  &  &  &  &  &  &  &  &  &  \\
\hline
 &  &  &  &  &  &  &  &  &  &  &  &  \\
\hline
December 2007~ &  &  &  &  &  &  &  &  &  &  &  &  \\
\hline
Sun~ &  & Mon~ &  & Tue~ &  & Wed~ &  & Thu~ &  & Fri~ &  & Sat~ \\
\hline
 &  &  &  &  &  &  &  &  &  &  &  & 1~ \\
\hline
 &  &  &  &  &  &  &  &  &  &  &  &  \\
\hline
2~ &  & 3~ &  & 4~ &  & 5~ &  & 6~ &  & 7~ &  & 8~ \\
\hline
 &  &  &  &  &  &  &  &  &  &  &  &  \\
\hline
9~ &  & 10~ &  & 11~ &  & 12~ &  & 13~ &  & 14~ &  & 15~ \\
\hline
 &  &  &  &  &  &  &  &  &  &  &  &  \\
\hline
16~ &  & 17~ &  & 18~ &  & 19~ &  & 20~ &  & 21~ &  & 22~ \\
\hline
 &  &  &  &  &  &  &  &  &  &  &  &  \\
\hline
23~ &  & 24~ &  & 25~ &  & 26~ &  & 27~ &  & 28~ &  & 29~ \\
\hline
 &  &  &  &  &  &  &  &  &  &  &  &  \\
\hline
30~ &  & 31~ &  &  &  &  &  &  &  &  &  &  \\
\hline
\end{longtable}
}

\section{Sheet6}

\begin{longtable}{|*{1}{>{\raggedright\arraybackslash}p{\dimexpr(\linewidth-2\tabcolsep-2\arrayrulewidth)/1\relax}|}}
\hline
Structural Elements~ \\
\hline
 \\
\hline
Miscellaneous structural elements you can add to your document, like footnotes, endnotes,~dropcaps~and the like.~~ \\
\hline
Footnotes \& Endnotes~ \\
\hline
\textcolor{hex000000}{{\fontsize{9pt}{10.8pt}\selectfont 1i}} \\
\hline
Dropcaps~ \\
\hline
D~ \\
\hline
rop~caps are used to emphasize the leading paragraph at the start of a section. In~Word~it is possible to specify how many lines of text a drop-cap should use.~ \\
\hline
Links~ \\
\hline
\textcolor{hex0000FF}{\uline{calibre download pageparagraph level formatting}} \\
\hline
\end{longtable}

\section{Sheet7}

\begin{longtable}{|*{1}{>{\raggedright\arraybackslash}p{\dimexpr(\linewidth-2\tabcolsep-2\arrayrulewidth)/1\relax}|}}
\hline
Images~ \\
\hline
Centered images like this are useful for large pictures that should be a focus of attention.~~ \\
\hline
 \\
\hline
% alt: 165 Inspirational Quotes To Keep You Motivated In Life
\includegraphics[width={\ifdim 473.25pt>\linewidth\linewidth\else 473.25pt\fi},height={\ifdim 473.25pt>0.9\textheight0.9\textheight\else 473.25pt\fi},keepaspectratio]{images/image1.png} \\
\hline
 \\
\hline
 \\
\hline
 \\
\hline
 \\
\hline
 \\
\hline
 \\
\hline
 \\
\hline
 \\
\hline
 \\
\hline
 \\
\hline
 \\
\hline
 \\
\hline
 \\
\hline
 \\
\hline
 \\
\hline
 \\
\hline
 \\
\hline
 \\
\hline
 \\
\hline
 \\
\hline
 \\
\hline
\textcolor{hex808080}{\textit{inserted}} \\
\hline
\end{longtable}

\section{Sheet8}

\begin{longtable}{|*{2}{>{\raggedright\arraybackslash}p{\dimexpr(\linewidth-4\tabcolsep-3\arrayrulewidth)/2\relax}|}}
\hline
Lists~ &  \\
\hline
All types of lists are supported by the conversion,~with the exception of~lists that use fancy~bullets,~these get converted to regular bullets.~ &  \\
\hline
Bulleted List~ &  \\
\hline
One~ &  \\
\hline
Two~ &  \\
\hline
Numbered List~ &  \\
\hline
1 & One, with~a very long~line to~demonstrate~that the hanging indent for the list is working correctly~ \\
\hline
2 & Two~ \\
\hline
Multi-level Lists~ &  \\
\hline
1 & One~ \\
\hline
1 & Two~ \\
\hline
1 & Three~ \\
\hline
2 & Four with~a very long~line to~demonstrate~that the hanging indent for the list is working correctly.~ \\
\hline
3 & Five~ \\
\hline
2 & Six~ \\
\hline
A~Multi-level~list with bullets:~ &  \\
\hline
One~ &  \\
\hline
Two~ &  \\
\hline
This bullet uses an image as the bullet item~ &  \\
\hline
Four~ &  \\
\hline
Five~ &  \\
\hline
Continued Lists~ &  \\
\hline
i & One~ \\
\hline
ii & Two~ \\
\hline
An interruption in our regularly scheduled listing, for this essential and~very relevant~public service announcement.~ &  \\
\hline
iii & We now resume our normal programming~ \\
\hline
iv & Four~ \\
\hline
\end{longtable}

\section{Chart\_Sheet}

{\small\setlength{\tabcolsep}{4pt}
\begin{longtable}{|*{6}{>{\raggedright\arraybackslash}p{\dimexpr(\linewidth-12\tabcolsep-7\arrayrulewidth)/6\relax}|}}
\hline
Charts &  &  &  &  &  \\
\hline
 &  &  & Series 1 & Series 2 & Series 3 \\
\hline
 &  & Category 1 & 4.3 & 2.4 & 2 \\
\hline
\textbf{Chart Title}

\begin{tabular}{|*{4}{>{\raggedright\arraybackslash}p{\dimexpr(\linewidth-8\tabcolsep-5\arrayrulewidth)/4\relax}|}}
\hline
 & Series D2 & Series E2 & Series F2 \\
\hline
Category 1 & 4.3 & 2.4 & 2 \\
\hline
Category 2 & 2.5 & 4.4 & 2 \\
\hline
Category 3 & 3.5 & 1.8 & 3 \\
\hline
Category 4 & 4.5 & 2.8 & 5 \\
\hline
\end{tabular} &  & Category 2 & 2.5 & 4.4 & 2 \\
\hline
 &  & Category 3 & 3.5 & 1.8 & 3 \\
\hline
 &  & Category 4 & 4.5 & 2.8 & 5 \\
\hline
\end{longtable}
}

\end{document}
